% from main.tex
%%%%%%%%%%  scan idx 389  |  folio 376  %%%%%%%%%%
\[
\begin{aligned}
  \bigl(f_*(F(n))\bigr)-\bigl(f_*(G(n))\bigr)
  ={}&\sum_i n_i\bigl((f_*(L_i(n)))-(f_*(M_i(n)))\\
      &\hspace{35mm}+(f_*(N_i(n)))\bigr).
\end{aligned}
\]
D'après 1.2,
pour \(n\) assez grand, les conditions suivantes seront remplies :
\[
  \forall k\geq1,\ \forall i,
  \quad R^k f_*(L_i(n))=0;
\]
\[
  \forall i,
  \quad f_*(L_i(n)),\ f_*(M_i(n)),\ f_*(N_i(n)),\ f_*(F(n)),\ f_*(G(n))
  \text{ sont localement libres de type fini.}
\]
Il en résulte que pour \(n\) assez grand, ces éléments sont tous dans \(A'\), et vérifient la relation écrite plus haut, et que les suites~:
\[
  0\longrightarrow f_*(L_i(n))\longrightarrow f_*(M_i(n))\longrightarrow f_*(N_i(n))\longrightarrow0
\]
sont exactes. Cela entraine que dans \(K^{\bullet}(S)\), \(f_*(F(n))\) et \(f_*(G(n))\) ont même classe.

\textbf{1.11.} Supposons donnée une équation de dépendance linéaire sur \(K^{\bullet}(S)\) entre \(1,\xi,\ldots,\xi^r\) :
\begin{equation}\tag{1.11.1}
  f^*(x_0)+f^*(x_1)\xi+\cdots+f^*(x_r)\xi^r=0.
\end{equation}
On peut écrire chacun des \(x_j\) comme différence des classes de deux faisceaux, soit : \(x_j=\cl(E_j)-\cl(E'_j)\), \(E_j\) et \(E'_j\) étant localement libres de type fini sur \(S\). Posons
\[
  F=\bigoplus_{j=0}^r f^*(E_j)(j),
  \qquad
  F'=\bigoplus_{j=0}^r f^*(E'_j)(j).
\]
L'égalité (1.11.1) s'écrit alors :
\[
  \cl_{K^{\bullet}(X)}(F)=\cl_{K^{\bullet}(X)}(F').
\]
D'après le lemme 1.10, on obtient pour \(n\) assez grand :
\[
  \cl_{K^{\bullet}(S)}\bigl(f_*(F(n))\bigr)=\cl_{K^{\bullet}(S)}\bigl(f_*(F'(n))\bigr).
\]
